\section{Common job settings}\label{sec-diabat-settings}
The principal diabatization jobs use the shared electronic-state package of Section~\ref{sec-statepack} and the fragment definition of Section~\ref{sec-diabat-frag}. The following options are used by several methods. The Job Control subsection for each method links back to these entries rather than repeating their definitions.

\subsection{Common keywords}
\subsubsection*{Keyword summary}
The following table provides a compact index of the options shared by the principal diabatization jobs. Each keyword links to its detailed description below.
\begin{keywordoverview}
\keyitem{key-common-diabat_wfn}{diabat\_wfn}{Save the constructed quasi-diabatic wave functions.}
\keyitem{key-common-partition}{partition}{Select L\"owdin or Mulliken population analysis.}
\keyitem{key-common-partial_diag}{partial\_diag}{Partially diagonalize the Hamiltonian within assigned state classes.}
\keyitem{key-common-adiabatic_hamiltonian}{adiabatic\_hamiltonian}{Supply the input adiabatic Hamiltonian in eV.}
\keyitem{key-common-outdir}{outdir}{Select the output directory.}
\keyitem{key-common-tmpdir}{tmpdir}{Select the temporary directory.}
\end{keywordoverview}
\subsubsection*{Detailed keyword descriptions}
The tables below give the data type, effective default behavior, accepted values, and practical use of each shared keyword.
\keyword
{key-common-diabat_wfn}{diabat\_wfn}{Write the resulting quasi-diabatic wave functions for further analysis.}{\OpOptional}{\OpLogic}{\OpTrue}{\OpTrue{} or \OpFalse}{Keep enabled when follow-up wave-function analysis is required.}{diabat\_wfn = .true.}
\keyword
{key-common-partition}{partition}{Choose the AO-based fragment population partition.}{\OpOptional}{\OpString}{\texttt{lowdin}}{\texttt{lowdin} or \texttt{mulliken}}{Choose according to the intended population analysis. L\"owdin is generally preferred.}{partition = "lowdin"}
\keyword
{key-common-partial_diag}{partial\_diag}{Perform partial Hamiltonian diagonalization among states assigned to the same class. For a calculation with exactly two target states, this step is skipped by default unless the keyword is explicitly supplied.}{\OpOptional}{\OpLogic}{Enabled for multistate jobs. Skipped automatically for two-state jobs when omitted}{\OpTrue{} or \OpFalse}{Normally omit this keyword. Set it explicitly only when the intermediate property-localized representation or an explicitly requested two-state partial diagonalization is required.}{partial\_diag = .true.}
\keyword
{key-common-adiabatic_hamiltonian}
{adiabatic\_hamiltonian}
{Specify the file containing a user-provided adiabatic Hamiltonian to be used in place of the default diagonal matrix of selected adiabatic-state energies. The file must use the numerical data format described in Appendix~\ref{app-numfile}. It must contain a real $N\times N$ matrix following the internal ordering of the $N$ adiabatic states being diabatized. \textbf{Every matrix element must be expressed in eV.} Supplying this file changes only the adiabatic Hamiltonian and does not modify the wave functions or state descriptors.}{\OpOptional}{\OpString}{Not used}{Readable numfile containing a real matrix}{Verify the file path, state order, matrix dimension, and energy units.}{adiabatic\_hamiltonian = "Hadia.dat"}
\keyword
{key-common-outdir}{outdir}{Directory for job-specific output files.}{\OpOptional}{\OpString}{Method name followed by .out}{Writable directory path}{Use a separate directory for each case.}{outdir = "fphd.out"}
\keyword
{key-common-tmpdir}{tmpdir}{Directory for temporary files.}{\OpOptional}{\OpString}{Method name followed by .tmp}{Writable directory path}{Use an appropriate high-capacity workspace.}{tmpdir = "fphd.tmp"}

\subsection{Restricting which states may mix}\label{sec-stateset}
The optional \texttt{\$\$stateset} block defines groups of states that may mix within a diabatization. Its indices are \emph{internal state IDs}, as explained in Section~\ref{sec-internal-ids}, not the external indices selected by \keyref{key-statepack-state_index}{state\_index}. Check the state-loading order reported on standard output before using the block.
\begin{lst-script}[escapechar=@]
@\scriptnum{script-stateset}@
$$stateset
  num_state_set = 2
  set 1 = 1..3
  set 2 = 4..6
$$
\end{lst-script}
The \texttt{\$\$stateset} block is normally unnecessary and should be used only when the intended calculation requires an explicit restriction on which states may mix. When it is used, verify the reported state groups after loading.
